% Original scenario reasoning; Dhaka 2025 printed figure checked separately.
\subsubsection{কাজ অনুযায়ী পেজের সংযোগ}
পেজের সংখ্যা দিয়ে কাঠামোর ধরন নির্ধারিত হয় না। কোন পেজ থেকে কোন পেজে সরাসরি যাওয়ার লিংক আছে এবং ব্যবহারকারী কোন কাজটি করবেন—এই দুটি বিষয় কাঠামো নির্বাচনের ভিত্তি।

\begin{multicols}{2}
\noindent\begin{minipage}[t]{\linewidth}\vspace{0pt}
\begin{keypointbox}{ধারাবাহিক কুইজের ক্ষেত্রে}
নির্দিষ্ট ক্রমে প্রশ্ন ১, ২, ৩ দেখাতে চাইলে লিনিয়ার কাঠামো উপযোগী। Next/Previous লিংক পরের বা আগের প্রশ্নে নিয়ে যায়; নির্ধারিত ধাপ বোঝা সহজ হয়। তবে শুধু পেজের এই সংযোগ দিয়ে সময়, নম্বর বা উত্তর যাচাই স্বয়ংক্রিয় হয় না; সেসবের জন্য অতিরিক্ত ব্যবস্থা দরকার।
\end{keypointbox}
\end{minipage}\par
\noindent\begin{minipage}[t]{\linewidth}\vspace{0pt}
\begin{keypointbox}{বিষয় বেছে অনুশীলনের ক্ষেত্রে}
Home থেকে বিষয়, তারপর সেই বিষয়ের প্রশ্নে যেতে হলে ট্রি কাঠামো স্বাভাবিক। ব্যবহারকারী নিজের বিষয় বেছে নিতে পারে। বিষয়ভিত্তিক ট্রির প্রতিটি প্রশ্নমালায় ধারাবাহিক Next/Previous যোগ করলে ট্রি ও লিনিয়ারের সমন্বয়ে হাইব্রিড কাঠামো হয়।
\end{keypointbox}
\end{minipage}\par
\end{multicols}
হোমপেজ থেকে সব পেজে যাওয়ার লিংক থাকলেও পেজগুলোর মধ্যে সরাসরি সংযোগ নাও থাকতে পারে। এমন অবস্থায় হোমপেজে ফেরার লিংক দিয়ে আগে সেখানে ফিরতে হয়, তারপর গন্তব্য পেজে যেতে হয়। সরাসরি আন্তঃপেজ লিংক থাকলে ধাপ কমে, তবে অতিরিক্ত লিংক পেজ জটিলও করতে পারে। কুইজে বিষয় বা কঠিনতার স্তর বাছাই করতে হলে ট্রি কাঠামো কাজে লাগে। শুধু পেজের সংযোগ দিয়ে সময়, নম্বর বা উত্তর যাচাই স্বয়ংক্রিয় হয় না; এসবের জন্য প্রোগ্রাম প্রয়োজন।
